\documentclass[10pt,letterpaper]{article}
\usepackage{amsmath,amssymb,amsthm}
\newtheorem{theorem}{Theorem}[section]
\newtheorem{corollary}[theorem]{Corollary}
\theoremstyle{remark}
\newtheorem*{remark}{Remark}
\newcommand{\Z}{\mathbb{Z}}

\title{Congruences and Fermat's Little Theorem}
\author{Daniel Okafor}
\date{June 18, 2018}

\begin{document}
\maketitle

\section{Modular arithmetic}
Two integers $a$ and $b$ are \emph{congruent modulo} $n$, written $a\equiv b\pmod n$, if $n$ divides $a-b$. Congruence is compatible with addition and multiplication, so we may compute in the ring $\Z/n\Z$.

Some elementary facts follow.
\begin{enumerate}
  \item If $a\equiv b \pmod n$ and $c\equiv d\pmod n$, then $a+c\equiv b+d\pmod n$.
  \item If $\gcd(a,n)=1$, there is an $a'$ with $aa'\equiv1\pmod n$.
  \item For a prime $p$, the ring $\Z/p\Z$ is a field:
    \begin{itemize}
      \item every nonzero element has an inverse;
      \item the multiplicative group has order $p-1$.
    \end{itemize}
\end{enumerate}

\section{Fermat's little theorem}
\begin{theorem}[Fermat]\label{thm:fermat}
Let $p$ be a prime and $a$ an integer not divisible by $p$. Then
\begin{equation}
  a^{p-1}\equiv 1 \pmod p. \tag{F}
\end{equation}
\end{theorem}

\begin{proof}
The numbers $a, 2a, \dots, (p-1)a$ are pairwise incongruent modulo $p$, hence a permutation of $1,2,\dots,p-1$. Multiplying them,
\[
  a^{p-1}\,(p-1)! \equiv (p-1)! \pmod p,
\]
and cancelling the unit $(p-1)!$ gives the claim.
\end{proof}

\begin{corollary}
For every integer $a$ and prime $p$, $a^p\equiv a\pmod p$.
\end{corollary}

\begin{remark}
The converse is false: the composite number $561=3\cdot11\cdot17$ satisfies $a^{560}\equiv1\pmod{561}$ for all $a$ coprime to it.
\end{remark}

\section{An example}
To find $7^{222}\bmod 11$ we note that $7^{10}\equiv1\pmod{11}$ by Theorem~\ref{thm:fermat}, and $222=22\cdot10+2$, so
\begin{align*}
  7^{222} &\equiv 7^2 \\
          &= 49 \equiv 5 \pmod{11}.
\end{align*}

\end{document}
